% !PRINTED 337
\weberpage{337}
We therefore set this determinant
\weq{2}{\Delta_{\lambda,\mu}(x) = Lf + M\Delta,}

wherein $L$ and $M$ are binary forms of the third degree in
$\lambda, \mu$.

If we apply an arbitrary linear transformation to the variables $x$,
then, on account of the invariant property of the function $\Delta$,
\weq{3}{\Delta_{\lambda,\mu}(y,b) = r^2\Delta_{\lambda,\mu\,r^2}(x,a),}

and from this it follows that the coefficients of the individual powers
and products of $\lambda, \mu$ are covariants of the original form, and
the coefficients $L, M$ of $f$ and $\Delta$ in these expressions, which
depend rationally on the coefficients $a$ of the form $f$, are therefore
\gesperrt{invariants} of this form. We have in this a means of actually computing
these invariants.

For if we denote by $a_i, A_i, B_i$ the coefficients of the
individual powers and products of the variables $x$ in $f, \Delta$ and
$\Delta_{\lambda,\mu}$, then from (2) there results a system of ten
equations:
\[
L a_i + M A_i = B_i,
\]
from which $L$ and $M$ can be determined. One obtains, if $i, k$ are two
distinct indices,
\weq{4}{\begin{aligned}
L\,(a_i A_k - a_k A_i) &= B_i A_k - B_k A_i\\
M\,(a_i A_k - a_k A_i) &= -B_i a_k + B_k a_i.
\end{aligned}}

From this one concludes further that $L, M$ are \gesperrt{integral}
functions of the coefficients $a_i$. For in (4) the right-hand sides are
integral functions, and if therefore $L$ or $M$, represented in the
simplest form, still had a denominator, then all the determinants
$a_iA_k - a_kA_i$, which are integral functions of the fourth degree in
the $a$, would have to have a common divisor. But that this is not
possible can be shown in any special case, e.g. for the form
\[
f = 3x_3(\alpha x_1^2 + \beta x_2^2) + a x_1^3 + b x_2^3,
\]
whose Hesse form is:
\[
-\beta^2(\alpha x_3 + a x_1)x_2^2 - \alpha^2(\beta x_3 + b x_2)x_1^2
\]
Here among the $a_iA_k - a_kA_i$ there occur, e.g., the two
functions $a^2\beta^2, b^2\alpha^2$, which have no common
divisor.
